\subsection*{能量层的轨道分解与振动周期}
\phantomsection\label{sec:energy-levels}
定理~\ref{thm:Hamilton函数沿轨道守恒} 说明轨道包含在能量层中；平衡点分类给出了临界点附近的形状。下面进一步区分能量层的各条轨道，并由沿轨道的速度计算周期。全局延拓的能量估计已见定理~\ref{thm:energy-global-extension}。
\begin{definition}
对一自由度系统 $\ddot x=-U'(x)$，记能量
$H(x,y)=y^2/2+U(x)$。若 $E$ 不等于任何 $U$ 的临界值，称为非临界能量。
满足 $U(a)=E$ 的位置称为转向点；当 $U'(a)\ne0$ 时称为简单转向点。
趋近平衡点的非平衡轨道称为分离轨道；两端趋于同一平衡点时称同宿轨道，
趋于两个不同平衡点时称异宿轨道。
\end{definition}
\begin{theorem}[非临界能量层]\label{thm:regular-energy-level}
设 $U\in C^2(\R)$，$E$ 为非临界能量。
集合 $\{x:U(x)<E\}$ 的每个有界连通分支 $(a,b)$ 对应一条闭轨，周期为
\[
 T(E)=2\int_a^b\frac{\dd x}{\sqrt{2(E-U(x))}}.
\]
无界分支对应直线型轨道；当该集合为整个实轴时，对应两条轨道
$y=\pm\sqrt{2(E-U(x))}$。
\end{theorem}
\begin{proof}
设 $(a,b)$ 是 $\{U<E\}$ 的一个有界连通分支。连续性给出
$U(a)=U(b)=E$；非临界假设又给出 $U'(a),U'(b)\ne0$。
由分支内部 $U<E$，必有 $U'(a)<0$、$U'(b)>0$。一阶展开于是给出
\[
 \begin{aligned}
 \lim_{x\downarrow a}\frac{E-U(x)}{x-a}&=-U'(a)>0,\\
 \lim_{x\uparrow b}\frac{E-U(x)}{b-x}&=U'(b)>0.
\end{aligned}
\]
故存在正常数 $c_1,c_2$，使靠近左端时
\[
 c_1(x-a)\leq E-U(x)\leq c_2(x-a);
\]
右端有对应估计。
因为 $\displaystyle\int_0^\delta s^{-1/2}\,ds=2\sqrt\delta<\infty$，下式收敛：
\[
 \tau=\int_a^b\frac{dx}{\sqrt{2(E-U(x))}}.
\]
在上半支 $y=\sqrt{2(E-U(x))}$ 上，$\dot x=y>0$，分离变量说明
从 $a$ 到 $b$ 的时间正是 $\tau$。到达 $(b,0)$ 时，向量场为
$(0,-U'(b))$，竖直向下且非零；存在唯一性把解接到下半支。
下半支满足 $\dot x=-\sqrt{2(E-U(x))}$，在同样的时间 $\tau$ 内返回 $(a,0)$。
在 $(a,0)$ 处向量场竖直向上，解遂重复。第一次回到初始状态需要完整走过两支，
故其最小正周期为 $2\tau$。

再看曲线的光滑性。能量层上
$\nabla H=(U'(x),y)$ 不为零，隐函数定理表明它在每一点附近都是可用单个坐标参数化的光滑曲线。
在转向点附近可用 $y$ 为坐标，因此两条显式平方根分支虽然在 $x$ 坐标中
斜率无穷，接合后的曲线仍光滑。两支加两个端点构成一条简单闭曲线。

若分支只有一个有限端点，正负两支经此端点接合，构成一条非闭合的直线型曲线。
若分支是整个实轴，则 $y$ 始终不为零，两支分别构成两条直线型曲线。
每一支上的非零向量场使任意两个有限位置之间的行程时间有限，
所以这里所述每条连通曲线确是一条最大轨道的像。
这些分支穷尽全部非临界能量层，分类因而完整。
\end{proof}
在临界能量层，$U'(a)=0$、$y=0$ 的点本身是一条平衡轨道，
不能同与之相接的分离轨道算成一条轨道。
唯一性说明非平衡运动不可能在有限时间到达平衡点；
若它在有限端点有此极限，则延拓后与常值解具有同一初值，矛盾。


图~\ref{fig:ode-tikz-energy-level} 将转向点与相平面中的上下两支对应起来；箭头表示真实的时间方向。
\begin{center}
\begin{minipage}{0.96\linewidth}
\centering
% TikZ mathematical illustration; compiled with the main document.
\begin{tikzpicture}[>=Stealth,font=\small]
\begin{scope}[x=1.2cm,y=1.45cm]
 \draw[->,gray] (-2,0)--(2.2,0) node[right]{$x$};
 \draw[->,gray] (0,-.18)--(0,2) node[above]{$U$};
 \draw[GreenMain,thick,domain=-1.8:1.8,samples=60] plot (\x,{\x*\x/2});
 \draw[GoldMain,thick] (-1.95,1)--(1.95,1) node[right]{$E$};
 \foreach \a/\lab in {-1.4142/a,1.4142/b}{
  \draw[dashed,gray] (\a,0)--(\a,1);
  \fill[GoldMain] (\a,1) circle(1.5pt);
  \node[below] at (\a,0){$\lab$};
 }
 \node at (0,-.55){允许位置：$a\leq x\leq b$};
\end{scope}
\begin{scope}[shift={(7,1.15)},scale=1.05]
 \draw[->,gray] (-1.85,0)--(1.95,0) node[right]{$x$};
 \draw[->,gray] (0,-1.8)--(0,1.9) node[above]{$y$};
 \draw[GreenMain,thick] (0,0) circle(1.4142);
 \draw[->,GreenMain,thick] (105:1.4142) arc[start angle=105,end angle=75,radius=1.4142];
 \draw[->,GreenMain,thick] (285:1.4142) arc[start angle=285,end angle=255,radius=1.4142];
 \fill[GoldMain] (-1.4142,0) circle(1.6pt) node[below left]{$(a,0)$};
 \fill[GoldMain] (1.4142,0) circle(1.6pt) node[below right]{$(b,0)$};
 \node at (0,-2.1){$y^2/2+U(x)=E$};
\end{scope}
\end{tikzpicture}

\captionof{figure}{势阱与闭合能量层：以 $U(x)=x^2/2$、$E=1$ 为例。}
\label{fig:ode-tikz-energy-level}
\end{minipage}
\end{center}

\begin{proposition}[数学摆的周期]
对 $\ddot x=-\sin x$，振幅 $0<A<\pi$ 时
\[
 T(A)=4\int_0^{\pi/2}
 \frac{\dd u}{\sqrt{1-k^2\sin^2u}},\qquad k=\sin(A/2).
\]
当 $A\to0$，
$T(A)=2\pi(1+A^2/16+O(A^4))$；当 $A\uparrow\pi$ 时周期按对数阶发散。
\end{proposition}
\begin{proof}
能量式给出 
\[
 \dot x^2=2(\cos x-\cos A).
\]
把四分之一周期的积分中令 $\sin(x/2)=k\sin u$，约去 $\cos u$，得到上述公式。
小 $k$ 时，积分号内一致展开为
$1+\frac12k^2\sin^2u+O(k^4)$，再用
$\displaystyle\int_0^{\pi/2}\sin^2u\,\dd u=\pi/4$ 和 $k^2=A^2/4+O(A^4)$ 即得渐近式。

令 $\delta=1-k^2$。在 $u=\pi/2$ 附近写 $s=\pi/2-u$，
分母平方是 $\delta+k^2\sin^2s$。
小 $s$ 时 $\sin s$ 与 $s$ 之比有正常数上下界，所以相应积分与
$\displaystyle\int_0^c(\delta+s^2)^{-1/2}\,\dd s$ 同阶。
后一积分为 $\operatorname{arsinh}(c/\sqrt\delta)$，是 $|\log\delta|$ 阶；
离开端点的部分有界，故证明对数发散。
\end{proof}
角度 $x$ 与 $x+2\pi$ 表示同一位置，因此摆的自然相空间是圆柱
$(\R/2\pi\Z)\times\R$。平面展开图中不闭合的旋转运动，在圆柱上可以闭合。


\subsection*{Lotka--Volterra系统的闭轨与平均值}
\phantomsection\label{sec:lv-energy}
回到绪论中的捕食者--食饵模型。虽然原坐标中的方程不直接具有标准 Hamilton 形式，仍可沿用“首次积分—紧等值线—周期运动”的论证。把参数记为正数 $a,b,c,d$，考虑
\[
 \dot x=x(-a+by),\qquad \dot y=y(c-dx),\qquad x,y>0.
\]
\begin{theorem}
函数
\[
 H(x,y)=dx-c\log x+by-a\log y
\]
是首次积分。正象限内除平衡点 $(c/d,a/b)$ 外的每条轨道均闭合。
任意周期轨道上，$x,y$ 的时间平均值分别为 $c/d,a/b$。
\end{theorem}
\begin{proof}
直接相乘可得
\[
 \dot H=(dx-c)(-a+by)+(by-a)(c-dx)=0.
\]
$H$ 的 Hessian 为 $\operatorname{diag}(c/x^2,a/y^2)$，正定；
$H$ 向象限边界或无穷远均趋于无穷，所以有唯一严格最小点 $(c/d,a/b)$。
每个非最小等值集紧，且梯度非零。
从最小点沿任意射线出发，$H$ 严格增加到无穷，因此每个较高等值集
在每个方向上恰交一点，并由隐函数定理构成一条光滑简单闭曲线。
向量场在其上非零且相切，紧性给出非零速度的正下界，所以有限时间内走完一周。

设周期为 $T$。对 $\dot x/x=-a+by$ 积分得
\[
 0=-aT+b\displaystyle\int_0^Ty(t)\,\dd t,
\]
对 $\dot y/y$ 积分得
\[
 0=cT-d\displaystyle\int_0^Tx(t)\,\dd t.
\]
除以 $T$ 即得平均值。
\end{proof}


图~\ref{fig:ode-tikz-lv-levels} 中闭轨围绕唯一极小点，各坐标的周期平均值均为 1。
\begin{center}
\begin{minipage}{0.96\linewidth}
\centering
% TikZ mathematical illustration; compiled with the main document.
\begin{tikzpicture}[>=Stealth,font=\small,x=1.8cm,y=1.55cm]
 \draw[->,gray] (0,0)--(3,0) node[right]{$x$};
 \draw[->,gray] (0,0)--(0,3) node[above]{$y$};
 \draw[dashed,GoldMain] (1,0)--(1,2.85);
 \draw[dashed,GoldMain] (0,1)--(2.85,1);
\draw[GreenMain,thick] plot coordinates {(1.45502,1.00000) (1.45462,1.01588) (1.45344,1.03171) (1.45152,1.04746) (1.44888,1.06309) (1.44557,1.07857) (1.44162,1.09387) (1.43706,1.10897) (1.43194,1.12386) (1.42628,1.13851) (1.42012,1.15291) (1.41349,1.16706) (1.40641,1.18095) (1.39891,1.19456) (1.39102,1.20791) (1.38274,1.22098) (1.37411,1.23377) (1.36514,1.24629) (1.35583,1.25852) (1.34620,1.27048) (1.33626,1.28216) (1.32603,1.29355) (1.31549,1.30467) (1.30467,1.31549) (1.29355,1.32603) (1.28216,1.33626) (1.27048,1.34620) (1.25852,1.35583) (1.24629,1.36514) (1.23377,1.37411) (1.22098,1.38274) (1.20791,1.39102) (1.19456,1.39891) (1.18095,1.40641) (1.16706,1.41349) (1.15291,1.42012) (1.13851,1.42628) (1.12386,1.43194) (1.10897,1.43706) (1.09387,1.44162) (1.07857,1.44557) (1.06309,1.44888) (1.04746,1.45152) (1.03171,1.45344) (1.01588,1.45462) (1.00000,1.45502) (0.98412,1.45461) (0.96830,1.45338) (0.95257,1.45129) (0.93699,1.44836) (0.92161,1.44457) (0.90649,1.43993) (0.89168,1.43446) (0.87722,1.42819) (0.86316,1.42115) (0.84954,1.41339) (0.83639,1.40495) (0.82374,1.39589) (0.81160,1.38627) (0.80000,1.37614) (0.78894,1.36556) (0.77842,1.35459) (0.76845,1.34329) (0.75900,1.33171) (0.75007,1.31989) (0.74165,1.30788) (0.73373,1.29573) (0.72627,1.28345) (0.71928,1.27109) (0.71272,1.25867) (0.70658,1.24621) (0.70085,1.23372) (0.69549,1.22124) (0.69051,1.20876) (0.68587,1.19629) (0.68158,1.18384) (0.67760,1.17142) (0.67394,1.15903) (0.67058,1.14667) (0.66751,1.13434) (0.66472,1.12203) (0.66221,1.10976) (0.65996,1.09750) (0.65798,1.08527) (0.65627,1.07306) (0.65481,1.06087) (0.65361,1.04868) (0.65268,1.03651) (0.65200,1.02433) (0.65160,1.01217) (0.65146,1.00000) (0.65160,0.98783) (0.65202,0.97567) (0.65274,0.96350) (0.65376,0.95134) (0.65509,0.93918) (0.65675,0.92704) (0.65875,0.91492) (0.66110,0.90282) (0.66383,0.89077) (0.66693,0.87877) (0.67044,0.86685) (0.67436,0.85502) (0.67872,0.84330) (0.68352,0.83172) (0.68877,0.82031) (0.69449,0.80910) (0.70069,0.79811) (0.70736,0.78738) (0.71451,0.77695) (0.72213,0.76684) (0.73021,0.75708) (0.73874,0.74771) (0.74771,0.73874) (0.75708,0.73021) (0.76684,0.72213) (0.77695,0.71451) (0.78738,0.70736) (0.79811,0.70069) (0.80910,0.69449) (0.82031,0.68877) (0.83172,0.68352) (0.84330,0.67872) (0.85502,0.67436) (0.86685,0.67044) (0.87877,0.66693) (0.89077,0.66383) (0.90282,0.66110) (0.91492,0.65875) (0.92704,0.65675) (0.93918,0.65509) (0.95134,0.65376) (0.96350,0.65274) (0.97567,0.65202) (0.98783,0.65160) (1.00000,0.65146) (1.01217,0.65160) (1.02433,0.65200) (1.03651,0.65268) (1.04868,0.65361) (1.06087,0.65481) (1.07306,0.65627) (1.08527,0.65798) (1.09750,0.65996) (1.10976,0.66221) (1.12203,0.66472) (1.13434,0.66751) (1.14667,0.67058) (1.15903,0.67394) (1.17142,0.67760) (1.18384,0.68158) (1.19629,0.68587) (1.20876,0.69051) (1.22124,0.69549) (1.23372,0.70085) (1.24621,0.70658) (1.25867,0.71272) (1.27109,0.71928) (1.28345,0.72627) (1.29573,0.73373) (1.30788,0.74165) (1.31989,0.75007) (1.33171,0.75900) (1.34329,0.76845) (1.35459,0.77842) (1.36556,0.78894) (1.37614,0.80000) (1.38627,0.81160) (1.39589,0.82374) (1.40495,0.83639) (1.41339,0.84954) (1.42115,0.86316) (1.42819,0.87722) (1.43446,0.89168) (1.43993,0.90649) (1.44457,0.92161) (1.44836,0.93699) (1.45129,0.95257) (1.45338,0.96830) (1.45461,0.98412) (1.45502,1.00000)};
\draw[->,GreenMain,thick] (1.28216,1.33626)--(1.38002,1.25069);
\draw[->,GreenMain,thick] (0.65481,1.06087)--(0.66888,1.19010);
\draw[->,GreenMain,thick] (1.12203,0.66472)--(0.99495,0.63732);
\draw[GreenMain,thick] plot coordinates {(1.98625,1.00000) (1.98508,1.03440) (1.98170,1.06865) (1.97629,1.10261) (1.96903,1.13619) (1.96010,1.16929) (1.94967,1.20186) (1.93789,1.23384) (1.92490,1.26521) (1.91083,1.29595) (1.89579,1.32604) (1.87988,1.35549) (1.86317,1.38431) (1.84574,1.41250) (1.82765,1.44007) (1.80895,1.46705) (1.78967,1.49344) (1.76986,1.51928) (1.74952,1.54456) (1.72869,1.56931) (1.70737,1.59355) (1.68556,1.61728) (1.66329,1.64053) (1.64053,1.66329) (1.61728,1.68556) (1.59355,1.70737) (1.56931,1.72869) (1.54456,1.74952) (1.51928,1.76986) (1.49344,1.78967) (1.46705,1.80895) (1.44007,1.82765) (1.41250,1.84574) (1.38431,1.86317) (1.35549,1.87988) (1.32604,1.89579) (1.29595,1.91083) (1.26521,1.92490) (1.23384,1.93789) (1.20186,1.94967) (1.16929,1.96010) (1.13619,1.96903) (1.10261,1.97629) (1.06865,1.98170) (1.03440,1.98508) (1.00000,1.98625) (0.96560,1.98503) (0.93138,1.98127) (0.89754,1.97486) (0.86428,1.96572) (0.83181,1.95383) (0.80036,1.93924) (0.77010,1.92207) (0.74122,1.90248) (0.71384,1.88071) (0.68806,1.85705) (0.66394,1.83179) (0.64148,1.80524) (0.62069,1.77771) (0.60149,1.74949) (0.58383,1.72083) (0.56761,1.69196) (0.55275,1.66307) (0.53915,1.63431) (0.52670,1.60580) (0.51531,1.57763) (0.50489,1.54987) (0.49536,1.52257) (0.48663,1.49575) (0.47864,1.46943) (0.47132,1.44361) (0.46461,1.41829) (0.45847,1.39345) (0.45284,1.36907) (0.44768,1.34513) (0.44296,1.32161) (0.43864,1.29848) (0.43471,1.27571) (0.43112,1.25328) (0.42788,1.23115) (0.42494,1.20930) (0.42232,1.18770) (0.41998,1.16632) (0.41792,1.14513) (0.41613,1.12410) (0.41462,1.10322) (0.41337,1.08244) (0.41240,1.06176) (0.41169,1.04114) (0.41126,1.02056) (0.41111,1.00000) (0.41126,0.97944) (0.41172,0.95886) (0.41251,0.93825) (0.41364,0.91759) (0.41514,0.89687) (0.41703,0.87609) (0.41936,0.85523) (0.42215,0.83431) (0.42546,0.81332) (0.42931,0.79229) (0.43378,0.77123) (0.43890,0.75018) (0.44475,0.72919) (0.45139,0.70830) (0.45888,0.68758) (0.46729,0.66712) (0.47667,0.64701) (0.48709,0.62735) (0.49857,0.60824) (0.51116,0.58981) (0.52484,0.57216) (0.53960,0.55539) (0.55539,0.53960) (0.57216,0.52484) (0.58981,0.51116) (0.60824,0.49857) (0.62735,0.48709) (0.64701,0.47667) (0.66712,0.46729) (0.68758,0.45888) (0.70830,0.45139) (0.72919,0.44475) (0.75018,0.43890) (0.77123,0.43378) (0.79229,0.42931) (0.81332,0.42546) (0.83431,0.42215) (0.85523,0.41936) (0.87609,0.41703) (0.89687,0.41514) (0.91759,0.41364) (0.93825,0.41251) (0.95886,0.41172) (0.97944,0.41126) (1.00000,0.41111) (1.02056,0.41126) (1.04114,0.41169) (1.06176,0.41240) (1.08244,0.41337) (1.10322,0.41462) (1.12410,0.41613) (1.14513,0.41792) (1.16632,0.41998) (1.18770,0.42232) (1.20930,0.42494) (1.23115,0.42788) (1.25328,0.43112) (1.27571,0.43471) (1.29848,0.43864) (1.32161,0.44296) (1.34513,0.44768) (1.36907,0.45284) (1.39345,0.45847) (1.41829,0.46461) (1.44361,0.47132) (1.46943,0.47864) (1.49575,0.48663) (1.52257,0.49536) (1.54987,0.50489) (1.57763,0.51531) (1.60580,0.52670) (1.63431,0.53915) (1.66307,0.55275) (1.69196,0.56761) (1.72083,0.58383) (1.74949,0.60149) (1.77771,0.62069) (1.80524,0.64148) (1.83179,0.66394) (1.85705,0.68806) (1.88071,0.71384) (1.90248,0.74122) (1.92207,0.77010) (1.93924,0.80036) (1.95383,0.83181) (1.96572,0.86428) (1.97486,0.89754) (1.98127,0.93138) (1.98503,0.96560) (1.98625,1.00000)};
\draw[->,GreenMain,thick] (1.59355,1.70737)--(1.69023,1.62045);
\draw[->,GreenMain,thick] (0.41462,1.10322)--(0.42322,1.23293);
\draw[->,GreenMain,thick] (1.20930,0.42494)--(1.08035,0.40845);
\draw[GreenMain,thick] plot coordinates {(2.60895,1.00000) (2.60649,1.05610) (2.59949,1.11185) (2.58853,1.16696) (2.57415,1.22123) (2.55681,1.27451) (2.53697,1.32669) (2.51497,1.37772) (2.49114,1.42758) (2.46575,1.47625) (2.43902,1.52376) (2.41112,1.57013) (2.38221,1.61540) (2.35241,1.65962) (2.32181,1.70282) (2.29049,1.74507) (2.25850,1.78640) (2.22589,1.82687) (2.19267,1.86653) (2.15887,1.90541) (2.12449,1.94356) (2.08953,1.98102) (2.05398,2.01782) (2.01782,2.05398) (1.98102,2.08953) (1.94356,2.12449) (1.90541,2.15887) (1.86653,2.19267) (1.82687,2.22589) (1.78640,2.25850) (1.74507,2.29049) (1.70282,2.32181) (1.65962,2.35241) (1.61540,2.38221) (1.57013,2.41112) (1.52376,2.43902) (1.47625,2.46575) (1.42758,2.49114) (1.37772,2.51497) (1.32669,2.53697) (1.27451,2.55681) (1.22123,2.57415) (1.16696,2.58853) (1.11185,2.59949) (1.05610,2.60649) (1.00000,2.60895) (0.94391,2.60630) (0.88826,2.59799) (0.83356,2.58356) (0.78038,2.56271) (0.72928,2.53535) (0.68081,2.50168) (0.63544,2.46216) (0.59352,2.41755) (0.55525,2.36880) (0.52067,2.31694) (0.48970,2.26305) (0.46213,2.20808) (0.43770,2.15289) (0.41611,2.09813) (0.39705,2.04433) (0.38023,1.99185) (0.36535,1.94091) (0.35218,1.89165) (0.34048,1.84414) (0.33008,1.79838) (0.32079,1.75434) (0.31248,1.71194) (0.30503,1.67112) (0.29833,1.63179) (0.29229,1.59384) (0.28683,1.55719) (0.28189,1.52173) (0.27742,1.48738) (0.27337,1.45405) (0.26969,1.42164) (0.26636,1.39009) (0.26333,1.35930) (0.26060,1.32920) (0.25813,1.29974) (0.25591,1.27083) (0.25392,1.24241) (0.25216,1.21444) (0.25061,1.18684) (0.24927,1.15957) (0.24813,1.13257) (0.24720,1.10580) (0.24646,1.07920) (0.24592,1.05273) (0.24560,1.02634) (0.24549,1.00000) (0.24560,0.97366) (0.24596,0.94727) (0.24657,0.92081) (0.24746,0.89424) (0.24865,0.86752) (0.25018,0.84062) (0.25208,0.81352) (0.25440,0.78620) (0.25719,0.75865) (0.26052,0.73085) (0.26446,0.70282) (0.26911,0.67458) (0.27456,0.64618) (0.28095,0.61767) (0.28841,0.58916) (0.29710,0.56078) (0.30718,0.53269) (0.31884,0.50511) (0.33223,0.47828) (0.34749,0.45248) (0.36472,0.42799) (0.38393,0.40507) (0.40507,0.38393) (0.42799,0.36472) (0.45248,0.34749) (0.47828,0.33223) (0.50511,0.31884) (0.53269,0.30718) (0.56078,0.29710) (0.58916,0.28841) (0.61767,0.28095) (0.64618,0.27456) (0.67458,0.26911) (0.70282,0.26446) (0.73085,0.26052) (0.75865,0.25719) (0.78620,0.25440) (0.81352,0.25208) (0.84062,0.25018) (0.86752,0.24865) (0.89424,0.24746) (0.92081,0.24657) (0.94727,0.24596) (0.97366,0.24560) (1.00000,0.24549) (1.02634,0.24560) (1.05273,0.24592) (1.07920,0.24646) (1.10580,0.24720) (1.13257,0.24813) (1.15957,0.24927) (1.18684,0.25061) (1.21444,0.25216) (1.24241,0.25392) (1.27083,0.25591) (1.29974,0.25813) (1.32920,0.26060) (1.35930,0.26333) (1.39009,0.26636) (1.42164,0.26969) (1.45405,0.27337) (1.48738,0.27742) (1.52173,0.28189) (1.55719,0.28683) (1.59384,0.29229) (1.63179,0.29833) (1.67112,0.30503) (1.71194,0.31248) (1.75434,0.32079) (1.79838,0.33008) (1.84414,0.34048) (1.89165,0.35218) (1.94091,0.36535) (1.99185,0.38023) (2.04433,0.39705) (2.09813,0.41611) (2.15289,0.43770) (2.20808,0.46213) (2.26305,0.48970) (2.31694,0.52067) (2.36880,0.55525) (2.41755,0.59352) (2.46216,0.63544) (2.50168,0.68081) (2.53535,0.72928) (2.56271,0.78038) (2.58356,0.83356) (2.59799,0.88826) (2.60630,0.94391) (2.60895,1.00000)};
\draw[->,GreenMain,thick] (1.94356,2.12449)--(2.03937,2.03662);
\draw[->,GreenMain,thick] (0.24813,1.13257)--(0.25315,1.26248);
\draw[->,GreenMain,thick] (1.27083,0.25591)--(1.14117,0.24641);

 \fill[GoldMain] (1,1) circle(1.8pt) node[above right]{$(1,1)$};
 \node[below] at (1,0){$1$};\node[left] at (0,1){$1$};
 \node[anchor=west] at (3.2,2.25){$H=x-\log x+y-\log y$};
 \node[anchor=west] at (3.2,1.65){$H=2.08,\ 2.30,\ 2.65$};
 \node[anchor=west] at (3.2,1.05){箭头沿系统运动方向};
\end{tikzpicture}

\captionof{figure}{参数 $a=b=c=d=1$ 时的 Lotka--Volterra 闭轨；各曲线为首次积分的等值线。}
\label{fig:ode-tikz-lv-levels}
\end{minipage}
\end{center}
